\begin{table}[t]
\centering
\scriptsize
\setlength{\tabcolsep}{3pt}
\resizebox{\columnwidth}{!}{%
\begin{tabular}{@{}lrrrr@{}}
\toprule
\textbf{Store} & \textbf{$G{=}1$} & \textbf{$G{=}32$} &
\textbf{$G{=}128$} & \textbf{$G{=}512$} \\
\midrule
32k, CPU-pinned & 8.9 & 9.2 & 10.9 & 94.0 \\
32k, GPU-resident & 8.4 & 7.7 & 5.5 & --- \\
128k, CPU-pinned & 27.6 & --- & --- & --- \\
128k, GPU-resident & 25.8 & --- & --- & --- \\
\bottomrule
\end{tabular}%
}
\caption{\textbf{Measured repeated-query break-even $Q^\star$.}
$G$ is generated tokens; a cell is the number of queries after which cumulative
CoMem time is lower than matched $j{=}0$ raw-text replay. Both arms share
selection, example, and pack; selection cancels in
$Q^\star=(W_{\mathrm{CoMem}}-I_{j=0})/
(T_{j=0}-T_{\mathrm{CoMem}})$. The 32k CPU cell uses
$W_{\mathrm{CoMem}}{=}2.253$\,s and a 2.16\,ms raw-token index build.
Values derive from three process medians and include store fetch, model Read,
and decode; raw component times and dispersion are archived. At 128k, only
$G{=}1$ is retained, so no generation-length trend is claimed there. ``---''
denotes an unavailable or non-finite retained cell.}
\label{tab:serving-crossover}
\end{table}
